% Extracto íntegro por residencia del cierre físico-matemático sucesor.
\section{Trou noir à double feuillet et purification avancée--retardée}
\label{sec:hmtf2-agujero-negro}

\subsection{État de feuillet et information}

Les canaux de feuillet sont
\[
 q_\pm=\exp(-A_{\rm rad}\mp C_{\rm rad}^*),\qquad
 p_\pm=\frac{q_\pm}{q_++q_-}.
\]
On définit
\[
 \rho_{\rm hoja}
 =p_+|+\rangle\langle+|+p_-|-\rangle\langle-|
\]
et la purification
\[
 |\Psi_{\rm ar}\rangle
 =\sqrt{p_+}|+\rangle|{\rm adv}\rangle
 +\mathrm e^{{\rm i}\chi_0}
  \sqrt{p_-}|-\rangle|{\rm ret}\rangle.
\]

\begin{proposition}[Purification exacte des \(2\) orientations temporelles]
\label{prop:hmtf2-purificacion-hojas}
L'état \(|\Psi_{\rm ar}\rangle\) est normalisé et sa réduction sur l'
orientation avancée--retardée est \(\rho_{\rm hoja}\). Son entropie est
\[
 S_{\rm hoja}
 =-\sum_\pm p_\pm\log p_\pm
 =\log(2\cosh C_{\rm rad}^*)
  -C_{\rm rad}^*\tanh C_{\rm rad}^*.
\]
La phase commune \(A_{\rm rad}\) s'annule et l'évolution globale réversible
possède un Landauer logique nul tant que l'étiquette de feuillet n'est pas effacée.
\end{proposition}

\begin{proof}
L'orthogonalité des deux orientations élimine les termes croisés lorsqu'on
prend la trace partielle. Substituer
\(p_\pm=\mathrm e^{\mp C_{\rm rad}^*}/(2\cosh C_{\rm rad}^*)\) donne l'
entropie.
\end{proof}

\subsection{Réalisation géométrique des \(2\) feuillets causaux}

Soit \(R(x)>0\) le rayon d'horizon publié par le sélecteur HMT--MD et
\[
 E_H(x)=\frac{c_{\rm int}^4R(x)}{2G_{\rm int}(x)}.
\]
Dans la région extérieure \(r>R(x)\), nous fixons
\[
 f(r)=1-\frac{R(x)}r,\qquad
 \mathrm ds^2=-f(r)c_{\rm int}^2\mathrm dt^2
 +f(r)^{-1}\mathrm dr^2+r^2\mathrm d\Omega_2^2.
\]
Soit
\[
 r_*=r+R\log\!\left(\frac rR-1\right)
\]
et soit \(\sigma(x)\in\{+1,-1\}\) le lecteur de feuillet. Nous définissons
\[
 U(x;t,r)
 =-\sigma(x)\exp\!\left[-\frac{c_{\rm int}t-r_*}{2R}\right],
 \qquad
 V(x;t,r)
 = \sigma(x)\exp\!\left[\frac{c_{\rm int}t+r_*}{2R}\right].
\]
Alors
\[
 -UV=\left(\frac rR-1\right)\exp(r/R).
\]

\begin{theorem}[Entrelaceur géométrique à double feuillet]
\label{thm:hmtf2-bh-doble-hoja}
La carte précédente réalise les deux feuillets HMT comme les deux régions extérieures
de l'extension de Kruskal. L'involution
\[
 C_{\rm BH}:(\sigma,U,V)\longmapsto(-\sigma,-U,-V)
\]
est isométrique, fixe \(r\), échange les deux régions extérieures et
préserve l'horizon \(UV=0\). La métrique étendue est
\[
 \mathrm ds^2
 =-\frac{4R^3}{r}\mathrm e^{-r/R}\,\mathrm dU\,\mathrm dV
 +r^2\mathrm d\Omega_2^2,
\]
régulière en \(r=R\), et satisfait les équations d'Einstein du vide hors
de la source. À l'horizon, elle conserve
\[
 E_H=\frac{c_{\rm int}^4R}{2G_{\rm int}},\qquad
 S_H=\frac{k_B\pi R^2}{\ell_P^2}.
\]
\end{theorem}

\begin{proof}
La formule de \(UV\) est la transformation standard du facteur
Schwarzschild et élimine la singularité de coordonnées en \(r=R\).
Le changement simultané de signe de \(U,V\) laisse \(UV\), \(r\) et la métrique
invariants. La substitution directe donne \(G_{\mu\nu}=0\) dans \(r>0\) hors de
la source. Les identités d'horizon dépendent seulement du rayon d'aire.
\end{proof}

\Needspace{25\baselineskip}
Soit \(\mathcal H_{\rm hoja}\) la somme directe des hilbertisations des
deux feuillets HMT et \(\mathcal H_{\rm ext}\) la somme directe des deux régions
extérieures avec la mesure transportée par la carte. Le changement de variables,
y compris son jacobien, induit une application unitaire
\[
 W_{\rm BH}:\mathcal H_{\rm hoja}\longrightarrow\mathcal H_{\rm ext}.
\]
Par conséquent,
\[
 \mathfrak I_{\rm BH}(A)=W_{\rm BH}AW_{\rm BH}^{\dagger}
\]
est une application unitale complètement positive, en fait un
\(*\)-isomorphisme, qui conserve l'état, le produit scalaire, l'involution de feuillet
et la dynamique modulaire lorsque l'état extérieur est transporté par \(W_{\rm BH}\).
Si les raffinements de feuillet et de carte conservent
\((\sigma,t,r,R)\), leurs isométries \(j_m^{\rm hoja}\) et
\(j_m^{\rm ext}\) satisfont
\[
 W_{{\rm BH},m+1}j_m^{\rm hoja}
 =j_m^{\rm ext}W_{{\rm BH},m}.
\]
La famille \(\mathfrak I_{{\rm BH},m}\) est alors naturelle sous
raffinement et transporte sans perte le même état enrichi.

Ainsi est achevée la construction mathématique
\[
\begin{aligned}
 \text{double feuillet HMT}
 &\longrightarrow
 \text{purification avancée--retardée}\\
 &\longrightarrow
 \text{deux extérieurs de Kruskal}
 \longrightarrow
 \text{horizon régulier et entrelaceur UCP}.
\end{aligned}
\]
L'identification à un objet astrophysique formé par effondrement, à son spectre
de modes quasi normaux, à son rayonnement, à ses échos ou à sa polarisation relève de la
\textsc{reconnaissance extérieure} et ne diminue pas la clôture géométrique.
